For the definite count, write \(R=\mathcal O_E\) and
let \(W\) denote the definite Hermitian space of rank \(k\) obtained
from the common target row. For a full integral lattice \(A\) in a
nondegenerate subspace, write \(A^\#\) for its Hermitian dual
in that subspace and \(\Delta(A)=|A^\#/A|\). In particular,
\(\Delta(0)=1\). Fix the genus \(\mathcal G\) of
Proposition~\ref{prop:flag-row-mass}, with \(T=T_u\). Thus
\[
     L^\#/L\simeq R/TR,\qquad D=\Delta(L)=T^{2f}
          \qquad(L\in\mathcal G),\qquad f=[F:\mathbf Q].
\]

An \emph{embedded datum of index \(x\)} over \(L\) consists of
a full integral lattice \(B\subset W\) of discriminant order \(D\)
and the inverse image \(\Gamma\subset B\times L\) of the graph
of a bijective isometry \(B/gB\to L/gL\) of the forms modulo \(g\).
Its diagonal lattice must satisfy
\[
 M_\Gamma=\{z\in W:(z,z)\in\Gamma\},\qquad
 [L:M_\Gamma]=x.
\]

The quotient \(L/M_\Gamma\) must be generated by at most \(s=m/2-1\)
elements over \(R\).
Write \(\mathcal T(L;x,x')\) for the set of triples consisting of
two such data \(d=(B,\Gamma)\), \(d'=(B',\Gamma')\), of indices
\(x,x'\), and a rational unitary isometry \(\eta:W\to W\) such that
\[
          \eta B'=B,\qquad (\eta\times1)\Gamma'=\Gamma .
\]

The lattices and isometries here are counted individually inside
the one target space. Normalization by target automorphisms occurs
only when averaging over the genus.

\begin{auxlemma}[Compatible embedded data in the definite genus]
\label{lem:zero-gram-definite-mean}
Let \(k\ge3\), \(g\ge2\), and \(x,x'\ge1\) be integers, and let
\(E/F\), \(T=T_u\) and \(\mathcal G\) be as in
Proposition~\ref{prop:flag-row-mass}. Put
\[
 \begin{gathered}
   c=1+\log_2(2f),\qquad v=\mathrm{rd}(E)/2\ge1,\qquad
   X=xx',\qquad Y=XD^2,\\
   M_b=\bigl(20f(2v)^2\bigr)^{2f(6+c)},\qquad
   \mathcal A=M_b(2v)^f,\\
   U_*=(32/g)^{kf}(4kf)^{4f}Y^7,\qquad
   V_*=(2/g)^{kf}(4kf)^{4f}DX .
 \end{gathered}
\]
For the fields of \eqref{eq:host-character-compositum},
the quadratic subfield \(\mathbf Q(\sqrt{-d})\) has discriminant
of absolute value \(d\equiv7\bmod16\) by
Proposition~\ref{prop:quadratic-prime-pair}. The discriminant tower formula
gives \(2v=\mathrm{rd}(E)\ge\sqrt d\).
For the probability proportional to \(1/|\operatorname{Aut}(L)|\)
on isometry classes in \(\mathcal G\),
\begin{equation}\label{eq:zero-gram-pair-mean}
   \big\langle |\mathcal T(L;x,x')|\big\rangle_{\mathcal G}
   \le X^{k/2}Y^{6s+10+4c}
      \left(1+\mathcal A\sum_{h=1}^{k-1}U_*^h+V_*^k\right).
\end{equation}
The bound applies when the two quotients have different minimal
numbers of generators, each at most \(s\).
\end{auxlemma}

\begin{proof}
\proofstep{Data with a prescribed diagonal lattice.}
Take \(M=M_\Gamma\) and \(u_0=2s+1\). The inclusions
\[
 M\subset L\subset L^\#\subset M^\#,\qquad
 |M^\#/M|=Dx^2,\qquad |M^\#/L|=Dx
\]
give \(M^\#/M\) at most \(u_0\) generators. Indeed the two outside
quotients are dual up to conjugation and the middle quotient is cyclic.
Every submodule of this finite module has the same generator bound
by elementary divisors locally and the Chinese remainder theorem.
Since \(M\subset B\subset M^\#\), a list of \(u_0\) elements
of \(M^\#/M\) suffices to specify \(B/M\). Once \(B\) is chosen,
fix lifts \(b_i\) of a generating list for \(B/M\). If \(l_i\)
represents the image of \(b_i\) modulo \(gL\), testing the isometry
on \(b_i\) and each \(z\in M\) gives
\[
 \langle b_i-l_i,z\rangle\in gR,\qquad
             (b_i-l_i)/g\pmod L\ \in M^\#/L .
\]

These classes determine the images of the \(b_i\), and the graph
already sends each \(z\in M\) to \(z\bmod gL\). The number of
possible data above \(M\) is consequently at most
\[
        (Dx^2)^{u_0}(Dx)^{u_0}\le(Dx)^{3u_0}.
\]

\proofstep{The common fixed space and its discriminant.}
Write \(M'=M_{\Gamma'}\), \(U=\ker(1-\eta)\), \(h=k-\dim_E U\).
On \(U\) the two graph conditions coincide. Put
\[
 M_0=M\cap U=M'\cap U,\qquad a=\Delta(L\cap U),\qquad
 y=[L\cap U:M_0]\le\min(x,x').
\]

Let \(P:W\to U^\perp\) be orthogonal projection. For \(z'\in M'\)
and \(z\in M\), the pairs \((\eta z',z')\) and \((z,z)\) belong
to \(\Gamma\). Comparing their pairings in its two coordinates gives
\[
    C(PM')\subseteq(PM)^\#,\qquad
                C=(1-\eta)|_{U^\perp}/g .
\]

Definiteness makes \(C\) invertible on \(U^\perp\). To take indices
also for these projected lattices, let \(\delta(A)\) be the real
covolume in the product of the \(f\) complex Hermitian spaces,
divided by \(v^{f\operatorname{rank}_E(EA)}\). This divides out
the covolume \(2^{-f}\sqrt{|\operatorname{disc}(E)|}=v^f\)
of \(R\) in one Hermitian line. The determinant formula gives
\[
  \delta(A^\#)=\delta(A)^{-1},\qquad
  \delta(M)=\delta(M_0)\delta(PM),\qquad
  \delta(M')=\delta(M_0)\delta(PM').
\]

The last two identities follow by projecting a real fundamental
parallelepiped along \(U\). For integral \(A\),
\(\delta(A)^2=\Delta(A)\). At each complex embedding, every
eigenvalue of \(\eta\) has absolute value one. With \(q_0=2/g\)
the projected inclusion therefore yields
\begin{equation}\label{eq:zero-gram-fixed-discriminant}
  ay^2=\delta(M_0)^2
   \le |N_{E/\mathbf Q}\det C|\,\delta(M)\delta(M')
   \le q_0^{2fh}DX .
\end{equation}

In particular, only integers \(1\le a\le DX\) can occur.

\proofstep{Triple count for a specified space \(U\).}
A rank-\(n\) projective \(R\)-lattice has at most \(I^{n+c}\)
sublattices of index \(I\), by the prime-by-prime chain count in
Proposition~\ref{prop:primitive-flag-means}. Thus for a fixed \(y\)
the number of \(M_0\) is at most
\[
   y^{k-h+c}\le (q_0^{2fh}DX/a)^{k/2}X^c.
\]
There are at most \(X\) choices for \(y\). Given \(M_0\), the image
of \(M\) in \(L/(L\cap U)\) has index \(x/y\), hence at most
\((x/y)^{h+c}\) possibilities. For a fixed image \(J\), lifts form
a torsor under \(\operatorname{Hom}_R(J,(L\cap U)/M_0)\).
Locally \(J\) is free of rank \(h\), so
\[
 |\operatorname{Hom}_R(J,(L\cap U)/M_0)|=y^h,\qquad
             (x/y)^{h+c} y^h\le x^{2h+c}.
\]
This bounds the choices for \(M\). Replacing \(x\) by \(x'\) gives
the corresponding bound for \(M'\).

With both data fixed, all compatible \(\eta\) that fix \(U\)
pointwise lie in one coset of the group of automorphisms of \(B\)
fixing \(U\) pointwise. This finite group acts faithfully on
\(B/(B\cap U)\), since orthogonality identifies its rational
action with the action on \(U^\perp\). For \(h>0\), put \(n=2fh\).
By Minkowski's theorem \cite[\S1.1]{serre-finite-orders}, its
order divides \((2n)!\). Indeed the bound on the exponent at
each prime \(p\) satisfies
\[
 \sum_{i\ge0}\left\lfloor\frac{n}{p^i(p-1)}\right\rfloor
       \le \sum_{i\ge1}\left\lfloor\frac{2n}{p^i}\right\rfloor
       =\operatorname{ord}_p((2n)!) .
\]
Consequently there are at most \((4hf)^{4hf}\) isometries, with
the value \(1\) when \(h=0\). Multiplying the bounds for the
intersection, its two lifts, the two graphs and the isometry gives
the following count over a specified \(U\):
\begin{equation}\label{eq:zero-gram-one-space}
  (q_0^{2fh}DX/a)^{k/2}
        X^{2h+2c+1}Y^{3u_0}(4hf)^{4hf}.
\end{equation}

\proofstep{The mean number of fixed spaces in intermediate dimension.}
For \(1\le h<k\), write \(z=k-h\), \(H_0=aD\) and
\(\lambda=\log_2(2f)=c-1\). Such spaces \(U\) of discriminant order
\(a\) correspond bijectively to their primitive lattices \(L\cap U\).
We need the constant in the gluing calculation of
Proposition~\ref{prop:primitive-flag-means}. Write \(m_i\) for
the standard unimodular mass in positive rank \(i\). The complementary
lattice \(L\cap U^\perp\) has discriminant order \(b\le H_0\):
the gluing group injects into both discriminant modules. In fact,
\[
    H=L/((L\cap U)\oplus(L\cap U^\perp)),\qquad
    ab=D|H|^2,\qquad |H|\le a .
\]

For the rational isometry type of \(U\), signatures are fixed and
at the finite places the determinant modulo norms classifies the
Hermitian space \cite{landherr}. At split places the class is trivial. At unramified
inert places it is determined by valuation parity and is trivial
away from primes over \(a\). At an inert prime \(\mathfrak P\)
an integral basis with Gram matrix \(H_{\mathfrak P}\) gives
\[
 \operatorname{ord}_{\mathfrak P}\det H_{\mathfrak P}
 =\operatorname{length}_{R_{\mathfrak P}}
     ((L\cap U)^\#/(L\cap U))_{\mathfrak P}.
\]

Thus these lengths also specify the determinant parities. The prime lists for those modules, of
total order \(a\), number at most \(a^\lambda\) by the index
count in Proposition~\ref{prop:primitive-flag-means}.
Local-to-global classification and Witt cancellation determine the
complementary rational space for each choice. Allow \(a^{1+\lambda}\)
choices. For each type the two fixed-discriminant mass bounds and
the sum over \(1\le b\le H_0\) cost at most
\[
          m_z a^{(z+\lambda)/2},\qquad
          H_0 m_h H_0^{(h+\lambda)/2}.
\]

To glue the two lattices, choose a subgroup of the direct sum of
their discriminant modules which injects into each factor. Its
projection into the rank-\(h\) discriminant module requires at most
\(h\) generators locally, hence globally. Even allowing
\((aH_0)^{2h+1}\) choices of subgroup suffices. Orbit--stabilizer
divides this count by the two automorphism orders before summing.
With \(N_{h,a}(L)=\#\{U:\dim U=z,\ \Delta(L\cap U)=a\}\)
the resulting unnormalized mass is at most
\[
 m_zm_h\,a^{z/2}H_0^{(9h+8+4\lambda)/2}.
\]

This calculation uses positive ranks \(z,h\), so includes \(h=1\)
and \(z=1\). By Proposition~\ref{prop:flag-row-mass},
the normalizing mass is
\[
     m_k\big((T^k-1)/(T-1)\big)^f\ge m_kD^{(k-1)/2}.
\]

Since \(a^{z/2} D^{-(k-1)/2}=(a/D)^{k/2}D^{1/2}a^{-h/2}\),
division gives
\[
 \langle N_{h,a}\rangle_{\mathcal G}
   \le \frac{m_zm_h}{m_k}(a/D)^{k/2}(aD)^{5h+5+2c}.
\]

It remains to control the rank dependence of the mass ratio. Put \(r=\min(z,h)\), \(t=k-r\) and \(R_i=m_i/m_{i-1}\)
for \(i\ge2\). Shimura's adjacent-rank formula as used in
Proposition~\ref{prop:primitive-flag-means} gives
\[
 R_i=\big[(i-1)!(v/\pi)^i\big]^f(2v)^{-f/2}
        L_F(i,\chi_{E/F}^i),\qquad
 2^{-f}\le L_F(i,\chi_{E/F}^i)\le2^f .
\]
The bounds follow from the Euler product and
\(\zeta(i)^f\le\zeta(2)^f<2^f\). The rank-one estimate there,
\(m_1\le2^fh_E\), and its Minkowski ideal count
\(h_E\le(2v)^{f(1+\lambda)}\) give \(m_1\le M_b\).
Using \(v\ge1\), factorial monotonicity, and \(r+t=k\),
\[
 \begin{aligned}
  \frac{m_zm_h}{m_k}
     &=\frac{m_1}{R_{t+1}}\prod_{i=2}^r\frac{R_i}{R_{t+i}}\\
     &\le m_1(2v)^f4^{rf}\pi^{(rt+1)f}
       \le \mathcal A\,16^{hkf}.
 \end{aligned}
\]
Indeed \(r(t+1)+1\le2hk\) and \(\pi<4\).

\proofstep{Summing the three ranges.}
For \(1\le h<k\), multiply the mean just obtained by
\eqref{eq:zero-gram-one-space}. The discriminant factors cancel
in the main exponent:
\[
  (q_0^{2fh}DX/a)^{k/2}(a/D)^{k/2}
         =q_0^{fhk}X^{k/2}.
\]
Use \(X\le Y\), \(aD\le Y\) and \(3u_0=6s+3\).
For each \(a\) the result is at most
\(X^{k/2}Y^{6s+9+4c}\mathcal A U_*^h\).
There are at most \(DX\le Y\) possible orders \(a\).

When \(h=0\), \(U=W\) and \(a=D\). Direct substitution in
\eqref{eq:zero-gram-one-space} gives at most
\(X^{k/2}Y^{6s+10+4c}\).
When \(h=k\), \(U=0\) is the sole space. Counting the two
diagonal lattices before imposing compatibility costs at most
\(X^{k+c}\). With the data and isometry bounds this gives
\[
       X^{k+c}Y^{3u_0}(4kf)^{4kf}.
\]
If any such triple occurs, \eqref{eq:zero-gram-fixed-discriminant}
with \(a=y=1\) gives \(1\le q_0^{2kf}DX\), hence
\(1\le X^{k/2}(q_0^{kf}D)^k\). Multiply by the latter factor
and collect powers to obtain at most
\(X^{k/2}Y^{6s+10+4c}V_*^k\). Adding the two endpoints and
the intermediate sum proves \eqref{eq:zero-gram-pair-mean}.
\end{proof}


We now attach the embedded data to actual pairs of flag terms.
Write \(s=m/2-1,\ d=\mathrm{rd}(E)\), and let \(\mathfrak d\)
be the different used for the integral frames. Fix a word \(u\).
For a nonunit term \(t\) with \(G_t=0\) and nonzero coefficient, let \(j\) be the last
nonzero leading minor of the mutual flag pairing and
\(\delta=(\delta_h)_{1\le h\le j}\) its list of leading minors.
In this calculation \(j\ge1\). Put
\[
 e_h=\sum_{a=1}^h(s+1-a),\quad
 A_h=(\delta_h)\mathfrak d^{-2e_h},\quad
 d_{v,h}=|\delta_h|_v d^{-2e_h},\quad
 D_\delta=N A_j=\prod_{v\mid\infty}d_{v,j}^2 .
\]
Here \(v\) runs over one complex embedding above each real place
and \(|\cdot|_v\) means ordinary complex modulus. The coefficient
ideals of the two frames show that the \(A_h\) are integral.
Omit zero minors when taking their norms. For \(H_1=b^{C_Hm}\) put
\[
 B_\delta=\prod_{\substack{1\le h\le j\\\delta_h\ne0}}N A_h,\qquad
 S_\delta=\prod_v\max\left(1,\frac{\max_{1\le h\le j}d_{v,h}}
                                      {(2H_1)^{2s}}\right),\qquad N=T_u^f .
\]

Write \(e_1,\ldots,e_s,e_{s+1}=z\) for the ordered vectors
of the first augmented flag and use primes for the second. Set
\(X_0=0\), \(X_i=\langle e_1,\ldots,e_i\rangle_E\) for \(i\le s\)
and \(X=X_s\). Representatives of the same graded frame satisfy
\[
 \widetilde e_i-e_i\in X_{i-1}\quad(1\le i\le s+1),
 \qquad \widetilde e'_i-e'_i\in X'_{i-1}.
\]

These changes preserve all leading determinants by unitriangularity.
Use Lemma~\ref{lem:zero-gram-positive-rank} to choose representatives
whose first \(j\) vectors have the signed block pairing and whose
remaining representatives satisfy
\[
 I=X\cap X'=\langle e_{j+1},\ldots,e_s\rangle_E,\qquad
 e'_i=\gamma e_i\ (j<i\le s+1),\qquad\gamma\bar\gamma=1.
\]

Indeed on the common tail the two vectors have this relation
modulo the preceding flag spaces; subtracting those components
gives the equalities in the chosen representatives. For the last
vector this uses \(z'-\gamma z\in X+X'\). The active vectors
and \(z\) span a nondegenerate space. In its orthogonal choose
duals of the tail, and subtract one half of their Gram matrix
in the tail directions to make the dual span \(I^*\) isotropic.
This last change concerns rational coordinates. The split reference lattice has the following summands. On
\(X\), and on the first \(j\) vectors of \(X'\), take the
coefficient ideals \(\mathfrak d^{-(s+1-a)}\). On the dual tail
take their inverses, and on \(z\) take \(\mathcal O_E\).
Call their sum \(P_0\). The full reference lattice is
\[
 L_{\rm ref}=P_0+
 \bigl((\Lambda+X)\cap(X+X'+I^*)^\perp\bigr).
\]

Let \(U\) be the common flag unipotent group and
\(\Gamma_A=U(F)\cap\operatorname{Stab}(A)\).
For these coordinates the inverse covolume, relative to the
reference lattice, is the index ratio
\[
 \frac{[\Gamma_{L_{\rm ref}}:\Gamma_{L_{\rm ref}}\cap\Gamma_\Lambda]}
      {[\Gamma_\Lambda:\Gamma_{L_{\rm ref}}\cap\Gamma_\Lambda]} .
\]

Write \(\omega_t\) for its infimum over adapted Witt coordinates
inducing the given graded vectors. An upper estimate for this
infimum may use any one adapted system.

\begin{auxlemma}[Flag terms above one embedded datum]
\label{lem:zero-gram-frame-count}
Fix \(u,\Lambda,j,\delta\) as above and one embedded datum
\((B,\Gamma)\) of index \(x\) in the definite space of \(\Lambda\).
Two second frames define the same class if a unipotent \(n\)
of the first flag preserves \(\Lambda\) and
\(n e'_i-e''_i\in X''_{i-1}\) for every \(1\le i\le s+1\). The sum of the weights \(\omega_t\)
of distinct frame classes above this datum is at most
\begin{equation}\label{eq:zero-gram-frame-bound}
 F_\delta(x)=
 (D_\delta/x)^{k/2}8^{s+1}f^2D_\delta^{3c+6}N^2
       (B_\delta D_\delta N)^{10000(s+1)^3},
 \qquad c=1+\log_2(2f).
\end{equation}
\end{auxlemma}

\begin{proof}
\proofstep{Indices in the common orthogonal quotient.}
In \(I^\perp/I\) let \(Q\) be the image of
\(\Lambda\cap I^\perp\). It is unimodular since
\(\Lambda\cap I\) is primitive in the unimodular lattice \(\Lambda\).
Write \(P,P'\) for the images of \(\Lambda\cap X,\Lambda\cap X'\),
\(Z=P_E+P'_E\), and \(K=Q\cap Z^\perp\).
Projection along \(P_E\) and \(P'_E\) gives the lattices
\[
 K_1=(Q+P_E)\cap Z^\perp,\quad
 K_2=(Q+P'_E)\cap Z^\perp,\quad
 B_0=[Q\cap Z:P+P'],\qquad p=[K_1:K].
\]

Set \(J=\Lambda\cap I\), \(\mathfrak a_i=\mathfrak d^{-(s+1-i)}\)
and \(A_i=\Lambda\cap X_i\). The index on the first side is
\[
 b_1=[A_s:A_j+J]=\prod_{i=j+1}^{s}[\mathfrak a_i:H_i],
 \qquad H_i=\{t\in E:\exists y\in J,\ y-te_i\in X_{i-1}\}.
\]

Indeed the successive factors between \(A_j+J\) and \(A_s+J=A_s\)
are obtained by taking the leading coefficient of \(A_i\) at \(i\).
Its image is \(\mathfrak a_i\) by the frame condition, and the
image of \(J\cap X_i\) is \(H_i\). For \(i>j\), the two filtrations
of \(I\) coincide and \(e'_i=\gamma e_i\), so the corresponding
ideal on the second side is \(H'_i=\gamma^{-1}H_i\). Thus
\[
 [\mathfrak a_i:H'_i]=[\mathfrak a_i:H_i],\qquad
 |N_{E/\mathbf Q}\gamma|=1 .
\]

This uses equality of covolumes; the two ideals need not be equal.
The product gives the same index \(b_1\) on \(X'\).

Let \(\bar A_j,\bar A'_j\) be the images of the two integral
active lattices in \(I^\perp/I\). Locally choose bases adapted
to the flag filtrations. Their leading coefficient ideals are
\(\mathfrak a_1,\ldots,\mathfrak a_j\) on each side. Consequently
the absolute norm of the cross-pairing determinant, with its two
coefficient ideals, is \(D_\delta\). The full Hermitian Gram matrix
has the cross pairing and its adjoint as its two off-diagonal blocks:
\[
 \sqrt{\Delta(\bar A_j+\bar A'_j)}=D_\delta,\qquad
 [P+P':\bar A_j+\bar A'_j]=b_1^2 .
\]

Both \(K_1,K_2\) are unimodular: projection along \(P_E\)
identifies \(K_1\) isometrically with the lattice \(P^\perp/P\)
in \(Q\), and likewise for \(P'\). Locally, complete a primitive
isotropic basis to hyperbolic pairs to see unimodularity of the
quotient. Therefore \([K_1:K]=[K_2:K]=|K^\#/K|^{1/2}\).
In the unimodular lattice \(Q\), the primitive perpendicular
lattices \(K\) and \(Q\cap Z\) have discriminant groups of equal
order: projection of \(Q/(K\oplus(Q\cap Z))\) identifies it
with each discriminant group, by perfect duality. On \(P+P'\)
the square root of this order is \(D_\delta/b_1^2\), by the
off-diagonal block determinant. Passing to \(Q\cap Z\) divides
that square root by its index \(B_0\). Thus
\[
          pB_0=D_\delta/b_1^2,\qquad [K_2:K]=p .
\]

For the intersection there is also an injection
\[
 (K_1\cap K_2)/K\ \longrightarrow\ (Q\cap Z)/(P+P').
\]

Indeed, subtract the two integral lifts of the same vector along
\(P_E,P'_E\). Changing either lift adds \(P\) or \(P'\). If the
difference vanishes modulo \(P+P'\), an integral common lift
lies in \(Z^\perp\).

Inside \(Z^\perp\) take the orthogonal complement of the
distinguished negative line, and denote it by \(W\).
Let \(\Theta=K\cap W\) and \(M=M_\Gamma\). Writing \(L\)
for the target definite lattice, we have
\[
 \Theta\subset M\subset (K_1\cap K_2)\cap W,\qquad L=K_1\cap W,\qquad
 a=[M:\Theta],\qquad xa=[L:\Theta]\le p,\quad a\le B_0 .
\]

For \(\Theta\) the same integral lift supplies both entries of
the congruence defining the graph. The inequality \(xa\le p\) is
the injection \(L/\Theta\hookrightarrow K_1/K\). Similarly
\(M/\Theta\hookrightarrow(K_1\cap K_2)/K\) gives \(a\le B_0\).
To count generators, write \(y\in K_1\) as \(y=q+v\) with
\(q\in Q\), \(v\in P_E\). Sending \(y\) to \(v\bmod P\)
embeds \(K_1/K\) in \(P_E/P\), since \(Q\cap P_E=P\).
Every finite subquotient of this rank-\(j\) torsion module is
generated by at most \(j\) elements at each prime, hence globally
by the Chinese remainder theorem. Apply this to \(L/M\), a
quotient of \(L/\Theta\). Multiplication of the two index inequalities gives
\begin{equation}\label{eq:zero-gram-effective-index}
                    x(a b_1)^2\le D_\delta .
\end{equation}

\proofstep{Counting integral realizations.}
Set \(D=D_\delta,\ B_*=B_\delta,\ n=N\mathfrak c\), with
\[
 \mathfrak c=(T_u)\prod_{h:\delta_h\ne0} A_h\overline {A_h},
 \qquad n=(NB_*)^2,\qquad \ell=2s+1,\qquad e=12s+50 .
\]

In the product \(1\le h\le j\). In simultaneous echelon form
the active pairing consists of signed antidiagonal blocks.
The nonzero minors give the block ends. For a block of length
\(h\) the first entry is a root of an equation of degree \(h\)
given by the ratio of successive end determinants, up to sign.
Thus compositions, signs and roots together cost at most \(8^s\)
(\(h\le2^{h-1}\)). The primitive terminal markings in the two middle
lattices give \(K_1\cap Ez=Rz\) and \(K_2\cap Ez=Rz'\), with
\(R=\mathcal O_E\) and \(z'=\gamma z\) in the common quotient.
Use the one lattice \(K\) to define two integral ideals by
\[
 \mathfrak a=\{a\in R:az\in K\},\qquad
 \mathfrak a'=\{a\in R:az'\in K\},\qquad
                   \mathfrak a=\gamma\mathfrak a'.
\]

The map \(a\mapsto az\bmod K\) embeds \(R/\mathfrak a\) in
\(K_1/K\), and \(a\mapsto az'\bmod K\) does the same for
\(R/\mathfrak a'\) in \(K_2/K\). Thus
\(N\mathfrak a,N\mathfrak a'\le p\le D\).
For each such pair of ideals the
generators differ by norm-one units, hence by roots of unity.
Using \(\varphi(t)\ge\sqrt{t/2}\) and \([E:\mathbf Q]=2f\) gives
\[
 \#\mu(E)\le8f^2,\qquad
 \#\{\gamma\}\le 8f^2D^{2(c+2)} .
\]

Here the number of ideals of a prescribed norm \(d_0\) is at
most \(d_0^c\): at a rational prime the exponent distributions
are bounded by the ordered choices of at most \(2f\) prime factors
at each step. The first row fixes its full middle lattice \(K_1\). On the other
side the datum gives \(B=K_2\cap W\), and \(\gamma\) fixes the
marked line lattice \(Rz'\). To recover the unimodular lattice
\(K_2\) containing \(Rz'\oplus B\), its gluing must be the graph
of an anti-isometry
\[
             (Rz')^\#/Rz'\simeq R/(T_u)\longrightarrow B^\#/B .
\]

Both summands are primitive, so both graph projections are
isomorphisms by unimodularity. The image of one generator determines
the map, giving at most \(|B^\#/B|=T_u^{2f}=N^2\) choices.

Fix these choices and \(a,b_1\). A section of the quotient onto
the projective lattice \(\Theta\) lifts the definite part of the
common orthogonal lattice. Changing the Witt complement by this
section is a common unipotent transformation: the rectangular block
is the difference of the lifts and its opposite is minus its adjoint.
The quadratic correction on the common tail can be taken with
a coefficient \(t\in\mathcal O_E\) satisfying \(t+\bar t=1\).
In coordinates with \(A\) the rectangular displacement, the
quadratic block \(C\) must solve
\[
 C+C^*=-A^* A,\qquad C=-t A^* A .
\]

Such \(t\) exists because \(E/F\) is unramified at finite places
and the integral trace is locally onto, including above 2.
In the resulting orthogonal decomposition the integral definite
intersection is \(\Theta\). Write \(P_{\rm int}\) for the
intersection with its complementary rational \(\ell\)-plane.
We record the denominator calculation in this plane:
\begin{equation}\label{eq:zero-gram-lattice-sandwich}
             \mathfrak c^{e}P_0\subset P_{\rm int}
                           \subset\mathfrak c^{-e}P_0 .
\end{equation}

For the active flags, triangular elimination uses the leading
block inverses. Each inverse is the adjugate divided by the
nonzero leading determinant, with exactly the indicated powers
of the different on its rows and columns. In particular, at a
block end \(h\) each oriented denominator satisfies
\[
 \frac{\mathfrak c\,\mathfrak d^{2e_h}}{\delta_h}
       \subset\mathcal O_E,\qquad
 \frac{\mathfrak c\,\mathfrak d^{2e_h}}{\overline{\delta_h}}
       \subset\mathcal O_E .
\]

One factor of
\(\mathfrak c\) clears the denominator of that pivot or its
conjugate. Applying the row and column eliminations at each flag
step costs at most two powers. Including the negative vector
and the common-tail frame costs two more. The two active lattices
and their central reference vectors therefore have denominator
exponent \(n_0=2s+2\).

For clarity, adjoin a dual tail vector \(y\) by choosing a lift
\(\Phi y\) with the prescribed pairings. First lift the pairings on \(\mathfrak c I^*\) into the
unimodular lattice. Then project to the complementary \(\ell\)-plane.
In terms of the two projections this lift has
\[
   \Phi=\operatorname{pr}_{\ell}\circ\Psi,\qquad
     \Psi(\mathfrak c I^*_{\rm int})\subseteq\Lambda .
\]

To check the denominator of this projection, take an integral
functional on \(\Theta\), lift it to \(\Lambda\) by primitivity,
and subtract its definite component. Modulo the common tail the
difference is supported on the two active flags and \(z\).
Solving for their coefficients uses the two mutual pairing matrices
and the norm of \(z\). Successive multiplication by the ideals of
the two oriented determinants and by \((T_u)\) clears them:
\[
 \mathfrak c^3\Theta^\#\subset\Theta,\qquad
     \mathfrak c^4\Phi(I^*_{\rm int})\subset\Lambda .
\]

Here \(I^*_{\rm int}\) has the dual frame ideals. The second
inclusion spends one power on the initial lift and three on
projecting it off the definite part. Put \(w_y=\Phi y-y\), and let \(q\) be the projection onto the
common tail dual to the projection onto the chosen dual tail.
Pairing with the central vectors costs \(n_0+4\) powers. Passing
from the paired lattice to \(P_0\) costs one, since
\(\mathfrak c P_0^\#\subset P_0\) on this plane. Consequently
the denominator exponents of the following two expressions are
bounded by
\[
    w_y-qw_y:\ n_0+5,\qquad
    (w_y,w_{y'}):\ 2(n_0+5)+1 .
\]

The common-tail components pair to zero with all the \(w_y\).
Subtracting the skew part of the pairing matrix by a central
shear leaves, in the pairing with \(y'\), the error
\(t((\Phi y,\Phi y')-(w_y,w_{y'}))\).
The corrected lift differs from \(y\) by a central vector with
denominator at most \(2n_0+12\). Expressing that vector in the
integral central lattice costs a further \(n_0\) powers.
In the reverse inclusion, pair with the corrected lifts and use
\(\mathfrak c P_0^\#\subset P_0\) once more. Both exponents are at
most \(6n_0+30\le e\). All these inclusions can be tested over
the local discrete valuation rings, which proves
\eqref{eq:zero-gram-lattice-sandwich} over \(\mathcal O_E\).

It follows that the choices of the two orthogonal intersections
and of the unimodular gluing are bounded respectively by
\[
       a^{k+c},\qquad n^{2e\ell^2},\qquad ((DN)^2)^\ell .
\]
For the first bound use elementary divisors in the rank \(k\)
lattice \(M\), summing over determinant ideals of norm \(a\).
For the second choose at most \(\ell\) generators in
\(\mathfrak c^{-e}P_0/\mathfrak c^eP_0\). For the third, a
unimodular gluing is the graph of an anti-isometry of discriminant
modules. Its source has at most \(\ell\) generators. Its target
has order \((xa)^2T_u^{2f}\le(DN)^2\).

\proofstep{Relative volumes.}
Filter the common unipotent group by its action on the span of the
two flags and \(z\), then by its action on the definite complement.
A shear \(u\) fixing the definite complement pointwise acts
only on the \(\ell\)-plane. For \(\lambda\in\Lambda\), write
\(\lambda=w+y\) with \(y\) in that plane. Pairing with
\(\mathfrak c^e P_0\subset\Lambda\) gives
\[
 y\in \mathfrak c^{-e}P_0^\#\subset\mathfrak c^{-e-1}P_0,\qquad
 (u-1)P_0\subset\mathfrak c^{3e+10}P_0
 \ \Longrightarrow\
 (u-1)y\in\mathfrak c^e P_0\subset\Lambda .
\]

Self-duality gives equality \(u\Lambda=\Lambda\). This applies
directly on the central skew subgroup. On the top quotient lift
the change of the two flag spans and \(z\) to their dual-tail
coordinates while fixing the definite complement. The dual-tail
linear block is the inverse adjoint on the common tail. The
component along \(z\) is determined by its pairing divided
by \(-T_u\), and its isotropic correction is obtained with
\(t+\bar t=1\). In particular, with
\(\mathfrak n=\mathfrak c^{3e+10}\), the initial depth gives
\[
    \mathfrak c^{3e+11}\subset (T_u)\mathfrak n,\qquad
    (u z-z)/T_u\in\mathfrak n P_0 .
\]

The remaining entries of the lifted difference lie in
\(\mathfrak n P_0\): taking the inverse adjoint preserves the
identity modulo the star-stable ideal \(\mathfrak n\), and all
terms of the quadratic correction contain a changed entry.
Thus the kernel of reduction on the top block also lifts to the
integral stabilizer. Each outer quotient acts on a module of rank
at most \(\ell\), so that reduction of its matrix entries gives
\[
  \text{relative index on each outer layer}
    \ \le\ |\mathcal O_E/\mathfrak c^{3e+11}|^{\ell^2}
    = n^{(3e+11)\ell^2}.
\]

The middle quotient has one additive row for each of the
\(s-j\) common-tail vectors, given by the definite component of
the change of its dual. A row integral for both its frame ideal
and \(\Theta\) lifts to a lattice-preserving shear, by adjoining
minus its adjoint and the integral correction \(t\) above.
In row \(i\) let \(h_i\) be the index of the common-tail leading
ideal in the full frame ideal. Filtration of \(\Lambda\cap X\)
by its ordered vectors gives
\[
 \prod_{i=j+1}^s h_i
   =[\Lambda\cap X:(\Lambda\cap X_{\le j})+(\Lambda\cap I)]
   =b_1 .
\]

A change of the coefficient ideal of index \(h_i\) costs \(h_i^k\)
in a row of rank \(k\). Restricting its definite entries from
\(L\) to \(\Theta\) costs \([L:\Theta]=xa\) in each row. Thus
\[
  \text{middle relative index}\ \le\ \left(\prod_i h_i^k\right)(xa)^{s-j}
          = b_1^k(xa)^{s-j}\le b_1^k D^s .
\]

Relative indices multiply along kernels and images. Multiplying
the three lattice counts by the three relative-index bounds gives
at most
\[
 (a b_1)^k a^c(B_*DN)^{10000(s+1)^3}.
\]

Indeed \(n\le(B_*DN)^2\); the exponent after substitution is at
most \(4e\ell^2+2\ell+s+4(3e+11)\ell^2\).
There are at most \(D^2\) possible pairs \(a,b_1\).
Use \eqref{eq:zero-gram-effective-index} for the factor with
exponent \(k\), use \(a^c\le D^c\), and multiply by the block,
scalar and line counts. This gives \eqref{eq:zero-gram-frame-bound}.
\end{proof}

\begin{auxlemma}[Matched zero-Gram terms with two prescribed labels]
\label{lem:zero-gram-label-pairs}
Let \(u\) be a selected word and fix two lists \(\delta,\delta'\)
of positive active ranks \(j,j'\le s\). Take any finite set
\(\mathcal P\) of pairs of nonunit terms with \(G_t=G_{t'}=0\)
in the same finite matrix block, with those respective labels.
In particular, the two terms have the same ordered endpoint rows
\(\Lambda,\Lambda'\) and the same finite operator coordinate.
Let \(w_\Lambda\) denote the reciprocal order of the finite
row automorphism image and \(M_u=\sum_\Lambda w_\Lambda\).
Uniformly on sufficiently late rows of the marked family,
\begin{equation}\label{eq:zero-gram-weight-sum}
 \begin{gathered}
 \sum_{(t,t')\in\mathcal P}
 w_{\Lambda}w_{\Lambda'}\,\omega_t\omega_{t'}
     \le M_u E_\delta E_{\delta'},\\
 E_\delta=D_\delta^{k/2}H_C^{Cs^4}S_\delta^{16k},
       \qquad H_C=b^{Cm^2f}.
 \end{gathered}
\end{equation}
The constant \(C\) is independent of both labels and the word.
\end{auxlemma}

\begin{proof}
\proofstep{From flag pairs to embedded isometries.}
Keep the first row fixed. First form the graph in the common
quotient \(I^\perp/I\) of a single pair of flags; its integral
lattice is \(Q\) from Lemma~\ref{lem:zero-gram-frame-count}.
Set \(S=P_E+Ez\), \(S'=P'_E+Ez'\), and
\(W_t=(S+S')^\perp\) in that quotient. Write
\(L_t=(Q+P_E)\cap W_t\) and \(B_t=(Q+P'_E)\cap W_t\). For
\(y'\in B_t\), \(y\in L_t\), the finite identification is the relation
\begin{equation}\label{eq:zero-gram-congruence-graph}
 (y',y)\in\Gamma_t
 \ \Longleftrightarrow\
 \begin{cases}
 y'-l'\in P'_E,\quad y-l\in P_E\\
 \text{for some }l',l\in Q,\quad l-l'\in gQ .
 \end{cases}
\end{equation}

Indeed these are exactly the conditions on integral representatives
of the two endpoint classes modulo their respective isotropic spaces.
Using \(Q\) gives the same graph as using \(\Lambda\) before passing
to \(I^\perp/I\): integral lifts of two congruent classes differ by
an element of \(g\Lambda+(\Lambda\cap I)\), and the last summand can
be subtracted on either side.

For the original augmented flags, set \(e'_{s+1}=z'\),
\(e_{s+1}=z\). At each \(1\le i\le s+1\) and \(t\in E\), the
marking says
\[
 \forall v'\in\Lambda\cap(te'_i+X'_{i-1})\quad
 \exists v\in\Lambda\cap(te_i+X_{i-1}):
                     \quad v-v'\in g\Lambda .
\] The
condition also works in reverse: lift the desired leading coefficient
using the common frame ideal, apply the forward condition and match
the difference in the preceding space by induction. In particular,
for \(v\in Q\cap S\) the resulting \(v'\in Q\cap S'\) satisfies
\(v'\equiv v\bmod gQ\). The same assertion holds for \(Q\cap P_E\)
and \(Q\cap P'_E\).

Take \(y'\in B_t\) and choose \(l'\in Q\) with \(y'-l'\in P'_E\).
Then \(l'\perp S'\) and
\[
 \langle v,l'\rangle=\langle v-v',l'\rangle\in gR
       \qquad(v\in Q\cap S).
\]
The primitive module \(Q\cap S\) is an \(R\)-direct summand of
\(Q\) since its quotient is projective over the Dedekind domain \(R\).
Extend the integral functional \(v\mapsto\langle v,l'\rangle/g\)
to \(Q\). Unimodularity gives a representative \(m_0\in Q\).
Consequently
\[
      l=l'-gm_0\in Q\cap S^\perp .
\]
The perfect rational pairing of \(P_E\) and \(P'_E\) gives a
unique \(p\in P_E\) with \(l-p\perp P'_E\). Since \(p\perp S\)
and \(z'=\gamma z\), the vector \(y=l-p\) belongs to \(W_t\).
It gives \((y',y)\in\Gamma_t\). Interchanging the flags proves
surjectivity of the second projection as well.

If \((0,y)\in\Gamma_t\), then \(l'\in Q\cap P'_E\).
Choose \(p\in Q\cap P_E\) congruent to \(l'\). Now
\[
      (l-p)/g\in Q,\qquad y/g-(l-p)/g\in P_E,\qquad y/g\in W_t ,
\]
so \(y\in gL_t\). The reversed argument gives the other kernel
\(gB_t\); conversely \(gB_t\times gL_t\subset\Gamma_t\) follows
directly from \eqref{eq:zero-gram-congruence-graph}. Pairings of
two representatives on one side agree with the pairings of their
lifts \(l\): their differences lie in the isotropic space \(P_E\)
orthogonal to the representatives. The analogous primed lifts are
congruent modulo \(gQ\), proving preservation of the form modulo \(gR\).

For different terms the rational spaces \(W_t\) may differ.
Projection modulo the original first isotropic space canonically
identifies them with the same definite quotient
\[
        W_0=(X+Ez)^\perp/X .
\]
It is injective since \(W_t\cap P_E=0\), and surjective by the
same rational correction \(l\mapsto l-p\). Both statements concern
Hermitian spaces; projection preserves the form. Under this map \(L_t\)
becomes the fixed target lattice \(L\), the image of
\(\Lambda\cap(X+Ez)^\perp\) in \(W_0\).
Projection carries \(B_t\) to an embedded lattice \(B\) and the
two coordinates of \(\Gamma_t\) to a graph \(\Gamma\). In each middle
lattice the distinguished line is primitive of norm \(-T_u\).
Unimodular gluing as in Proposition~\ref{prop:flag-row-mass} gives
\[
 B^\#/B\simeq L^\#/L\simeq R/(T_u),\qquad
 [L:M_\Gamma]=x,\qquad L/M_\Gamma\text{ has at most }j\text{ generators}.
\]
The two index assertions follow from
Lemma~\ref{lem:zero-gram-frame-count} and its proof. Apply the
same construction separately to the other term, with index \(x'\)
and at most \(j'\) generators.

The rational isometries from the second definite row to these
two embedded lattices, say \(h_t,h_{t'}\), give
\[
   \eta=h_t h_{t'}^{-1},\qquad
   \eta B'=B,\qquad(\eta\times1)\Gamma'=\Gamma .
\]
In these equations \(h_t\) carries the fixed definite quotient
of the second row to \(B\): transport an integral representative by
the rational isometry of the term, quotient by the transported
isotropic flag, and use \(W_t\simeq W_0\). The finite operator coordinate determines the correspondence
on the two middle lattices modulo \(g\), hence on their marked
definite parts. In the ordering (source, target) this last
correspondence is given by
\[
 (v\bmod g,w\bmod g)\text{ correspond}\ \Longleftrightarrow\
 (h_t v,w)\in\Gamma
\]
by \eqref{eq:zero-gram-congruence-graph}. This proves the graph identity for \(h_th_{t'}^{-1}\).

To bound the fibre, fix a target row and one class of second
graded frames. Two realizations differ by a unipotent preserving
the target lattice, followed by a change of representatives in the
graded frame. Transporting these maps back identifies the source
lattices with their graded vectors. To identify the actual marked
vectors, split the primitive isotropic lattice \(A\) with its
hyperbolic dual as in Proposition~\ref{prop:flag-row-mass}, writing
the integral lattice as \(A\oplus H\oplus A^\vee\). In a marked
source row \(A=\bigoplus_i\mathfrak a_i e_i\). The two choices
of representatives, with these same coefficient ideals, differ on
\(A\) by an integral unitriangular automorphism \(T\); its inverse
is integral by triangular elimination. Extend it by
\[
            T\oplus 1_H\oplus (T^*)^{-1}.
\]

The remaining difference of the two marked terminal vectors lies
in \(A\), say \(a_0\). Their common image \(z_H\) is primitive
in \(H\), so there is an integral map \(C:H\to A\) with
\(Cz_H=a_0\) (split the primitive summand \(Rz_H\)). Its unitary
extension, for \(\tau+\bar\tau=1\) in \(R\), is
\[
 (a,h,y)\longmapsto
 (a+Ch-\tau CC^*y,\ h-C^*y,\ y).
\]

The two linear cross terms cancel, and the remaining quadratic
coefficient is \(1-\tau-\bar\tau=0\). Thus the source marking
can be aligned by an integral first-flag unipotent.
The two sources are the same row \(\Lambda'\) in the chosen set
of representatives. Fix one realization. For any
other, the residual source automorphism fixes the graded frame.
On \(X^\perp/X\) it fixes the negative vector and acts by an
isometry on its definite complement. Two such automorphisms with the
same induced isometry differ by a first-flag unipotent: they agree
on every graded line, the middle quotient and the dual graded lines
by unitarity. Consequently, denoting a second graded-frame class by \([F']\),
\[
 \#\{\text{terms over }\Lambda,[F']\}
        \le|\Delta_{\Lambda'}|=1/w_{\Lambda'} .
\]

Here \(\Delta_{\Lambda'}\) is the finite row automorphism image.

For a fixed pair of frame classes and a fixed triple, choice of
the first term specifies \(h_t\); the equation
\(h_{t'}=\eta^{-1}h_t\) now fixes the second term by the same
double-coset argument. Thus the weighted fibre has total source
weight at most \(w_{\Lambda'}(1/w_{\Lambda'})=1\).
The two Hilbert--Schmidt row factors are
\[
  (\sqrt{w_\Lambda w_{\Lambda'}})^2=w_\Lambda w_{\Lambda'} .
\]
The second cancels the fibre bound; the first remains for the
genus average. The relative index has the same value on
integrally equivalent graded frames, by conjugation.
Lemma~\ref{lem:zero-gram-frame-count} bounds each frame-class
sum. Proposition~\ref{prop:flag-row-mass} identifies the
normalized average on first rows with that on their definite genus.
Thus division of the left side of \eqref{eq:zero-gram-weight-sum}
by \(M_u\) leaves at most
\[
 \sum_{1\le x\le D_\delta}\sum_{1\le x'\le D_{\delta'}}
  F_\delta(x)F_{\delta'}(x')\,
         \langle|\mathcal T(L;x,x')|\rangle_{\mathcal G}.
\]

\proofstep{Absorbing the two index sums.}
Take \(C_a\ge\max(2(C_H+1),C_T,1)\), where
\(T_u\le b^{C_Tm}\). The bounds for \(m,k,f\) in Lemma~\ref{lem:rank-window}
give \(m^4/f\ll x^{-33\epsilon/4}\),
\(m^5 f/k\ll x^{-37\epsilon/16}\), and
\((\log f)^2/s\ll(\log x)^2x^{-1/2+\epsilon/16}\).
Together with the marked level choice in
Paragraph~\ref{par:window-successive-bounds} they give the late range
\[
 \begin{gathered}
 m\ge42C_a,\quad m^4\le f,\quad m^5 f\le k,\quad
 42 C_a m^2\log b\le k,\\
 \log f\le\sqrt{s},\quad 2\le d\le b,\quad (kb)^8\le g .
 \end{gathered}
\]

Set \(A=b^{C_am^2f},\ U=A S_\delta^2,\ U'=A S_{\delta'}^2,\)
and \(Z=S_\delta S_{\delta'}\).
Taking the product of the local maxima gives
\[
 D_\delta\le U,\qquad B_\delta\le U^s,\qquad N\le A,\qquad
 F_\delta(x)\le (D_\delta/x)^{k/2}U^{250000s^4}.
\]

For the last inequality substitute the first three in
\eqref{eq:zero-gram-frame-bound}; \(c\le4s\) and
\(8^{s+1},f^2\le A\). In the genus estimate
\eqref{eq:zero-gram-pair-mean}, \(Y=xx'N^4\le A^6Z^2\).
The coefficient of \(Y^7\) obeys
\[
          (32/g)^{kf}(4kf)^{4f} A^{42}\le 2^{-f}.
\]

Indeed take the \(f\)-th root, use \(b^{42C_am^2}\le e^k\),
\(f\le k,\ b\ge2\) and \(g\ge(kb)^8\); the result is at most
\((32e/(2k)^8)^k(4k^2)^4\le1/2\) for \(k\ge9\).

It follows that \(U_*,V_*\le Z^{14}/2\) in
\eqref{eq:zero-gram-pair-mean}. Since \(Z\ge1\), the truncated
sum of their powers is bounded after extracting \(Z^{14k}\)
by a geometric series of ratio \(1/2\). The two other factors obey
\[
 Y^{6s+10+4c}\le(UU')^{96s},\qquad
 2+2(20fd^2)^{2f(6+c)}d^f\le A^{183}.
\]
For the first use \(Y\le(UU')^3,\ c\le4s\); for the second
use \(f\le b^{2s},\ d\le b,\ 20\le b^5\), and \(m=2s+2\).
The genus bound contains \((xx')^{k/2}\), cancelling precisely
the denominator in the two \(F\)'s. At most \(UU'\) index
pairs occur. Hence the contribution per label is at most
\[
 D_\delta^{k/2}U^{300000s^4}S_\delta^{14k},
 \qquad 1+96s+183+250000s^4\le300000s^4 .
\]
Finally \(m^4\le f,\ m^5 f\le k,\ m\ge42\) imply
\(300000s^4\le k\). With \(C=300000C_a\) the last expression
is at most \(E_\delta\), and the same calculation applies
to \(\delta'\).
\end{proof}
